\appendix
\section{Códigos de problema 1 Inciso 2} \label{ap: sec1prob2}
Se utilizó Python y se usó la IA mayormente para resolver dudas y graficar de mejor manera. Todo el código y lógica fue hecho por los autores y estos se hacen responsables de su veracidad.

\begin{lstlisting}[language=Python, caption={Problema 1 Inciso 2.}]
import numpy as np
from scipy.optimize import brentq
import matplotlib.pyplot as plt
from matplotlib.ticker import MultipleLocator

def energias_permisibles(mumax: float,*,eps=1e-12,a=1,hbar=1,M=1/8 * np.pi**2) -> np.ndarray:
    '''
    Se calcula las energías permisibles para el pozo finito. Se sabe que la función tan(pi/2 * x),
    tiene polos en los valores pares de x y siempre es positivo en el intervalo 2n, 2n+1.

    Args:
        mumax:      the maximum value of mumax. (sqrt(U0/V0))
        eps:        a small tolerance parameter
        a:          width of the finite well
        hbar:       the value of Planck reduced constant
        M:          the mass of the particle
    returns:
        Tuple(array,array) with even and odd energies respectively.

    '''
    f = lambda x: x*np.sin(np.pi*x/2) - np.sqrt(mumax**2 - x**2)*np.cos(np.pi*x/2)
    g = lambda x:  np.sqrt(mumax**2 - x**2)*np.sin(np.pi*x/2) +x*np.cos(np.pi*x/2)
    raicespar = []
    raicesimpar = []
    n = 0
    while 2*n < mumax:
        #estados pares
        low = 2*n + 1e-12
        high = min(2*n + 1-eps, mumax)
        if f(low)*f(high) < 0:
            raicespar.append(brentq(f, low, high))
        #estados impares
        low+=1; high= min(high+1,mumax)
        if low>mumax:
            n+=1
            break
        if g(low)*g(high) < 0:
            raicesimpar.append(brentq(g,low,high))
        n += 1
    E1 = (hbar*np.pi)**2/(2*M*(2*a)**2)    
    x_par, x_imp = np.array(raicespar),np.array(raicesimpar)
    energias = (E1*x_par**2, E1*x_imp**2)
    return energias





def plot_energies(energias,*,a=1,ax=None): 
    '''
    Plots allowed energie levels in horizontal lines.
    
    Args:
        energias:       array(s) that contain the energies
        a:              width of the well
        ax:             if included, the axis for graph
    Raises:
        Value error if energies is empty.
    '''
    energias =  np.sort(np.concatenate(energias))
    if len(energias) ==0:
        raise ValueError("Energías must at least have 1 value.")
    if ax is None:
        _, ax = plt.subplots()

    for (index,e) in enumerate(energias,start=1):
        ax.text(1, e, f"  n={index}", va="center_baseline")
    for e in energias:
        plt.hlines(e,-a,a,lw=0.6)
    plt.ylabel("Energía") 
    ax = graph_well(ax)
    plt.title("Valor de las Energías permisibles")
    plt.show()

def graph_well(ax,*,a=1,U0=110):
    def potential(x: np.ndarray, *, a: float, U0: float) -> np.ndarray:
        """
        Helper function that returns 0 for |x| <= a and U0 otherwise.

        Args:
            x: Positions at which to evaluate the potential.
            a: Half-width of the well.
            U0: Potential height outside the well.

        Returns:
            Array of potential values with the same shape as x.
        """
        return np.where(np.abs(x) <= a, 0.0, U0)
    steps = 2001
    x = np.linspace(-2 * a, 2 * a, steps)
    y = potential(x, a=a, U0=U0)

    ax.set_xlim(-2 * a, 2 * a)
    ax.set_ylim(-0.1 * U0, 1.3 * U0)
    ax.set_xlabel(r"$x$")
    ax.set_ylabel(r"$U(x)$")
    ax.grid(alpha=0.3)
    ax.plot(x,y,lw=1)
    return ax

def plot_eigenstates(energias,M,hbar,U0,a,*, n=8):
    '''
    Function that plots eigenstates in a 4x2 grid. To graph, it also calculates the eigenstates.

    Args:
        energias:   tuple that contains even and odd energies in arrays.
        M:          mass of the particle
        hbar:       value of Planck reduced constant
        U0:         value of the potential outside the well
        a:          width of the well
        n:          number of eigenstates to graph (max 8)

    Raises:
        ValueError if there is not enough energies or n>8
    ''' 

    if (len(energias[1]) < n/2):
        raise ValueError("No hay suficientes energías.")

    if (n>8):
        raise ValueError(f"El valor de $n$ es muy grande. El máximo es 8")
    knpar = np.sqrt(2*M*energias[0])/hbar
    qnpar = np.sqrt(2*M*(U0-energias[0]))/hbar
    knimpar = np.sqrt(2*M*energias[1])/hbar
    qnimpar = np.sqrt(2*M*(U0-energias[1]))/hbar
    Anpar = np.sqrt(qnpar/(a*qnpar+1))
    Bnpar = np.cos(knpar*a)/np.exp(-qnpar*a)*Anpar
    Animpar = np.sqrt(knimpar/(a*knimpar-np.tan(a*knimpar)))
    Bnimpar = np.sin(knimpar*a) / np.exp(-qnimpar*a) * Animpar
    steps= 2001
    fig, axes = plt.subplots(4, 2, figsize=(10, 12), sharex=True)
    x1 = np.linspace(-2*a,-a,steps)
    x2 = np.linspace(-a,a,steps)
    x3 = np.linspace(a,2*a,steps)
    count = 0
    for count,ax in zip(range(n),axes.flat):
        i = count // 2
        if count % 2 == 0:
            # ESTADO PAR
            psi = np.concatenate([Bnpar[i]*np.exp(qnpar[i]*x1),
                                  Anpar[i]*np.cos(knpar[i]*x2),
                                  Bnpar[i]*np.exp(-qnpar[i]*x3)])
        else:
            # ESTADO IMPAR
            psi = np.concatenate([-Bnimpar[i]*np.exp(qnimpar[i]*x1),
                                  Animpar[i]*np.sin(knimpar[i]*x2),
                                  Bnimpar[i]*np.exp(-qnimpar[i]*x3)])
        ax.plot(np.concatenate([x1, x2, x3]), psi, label=f"n={count+1}")
        ax.set_title(f"n={count+1}")
        # well line
        ax.axvline(-a, color="red", lw=0.4, ls="--")
        ax.axvline(a, color="red", lw=0.4, ls="--")
    fig.suptitle(f"Primeros {n} Eigenestados del pozo finito", fontsize=14)
    fig.tight_layout()
    for ax in axes[-1]:
        ax.set_xlabel("x")
    for ax in axes[:, 0]:
        ax.set_ylabel(r"$\psi_n(x)$")

    plt.legend()
    plt.show()


def eigenestados(
    energias: tuple[np.ndarray, np.ndarray],
    M: float,
    hbar: float,
    U0: float,
    a: float,
    *,
    n: int = 8,
    x_ext: float = 3,
    steps: int = 6001,
) -> tuple[np.ndarray, np.ndarray]:
    """
    Computes eigenstate for a finite potential well and returns them in an increasing energy order.

    Args:
        energias:   tuple that contains even and odd energies in arrays.
        M:          mass of the particle
        hbar:       value of Planck reduced constant
        U0:         value of the potential outside the well
        a:          width of the well
        n:          number of eigenstates to graph (max 8)
        x_ext:      the máximum value of exterior that the eigenstates takes
        steps:      the total amount of points evaluated between -x_ext and x_ext

    Returns:
        A tuple x,psi where x is an array of length steps and psi are the eigenstates in a shape (n,steps).

    Raises:
        ValueError: If there are not enough energies for n states, or if
            any energy lies outside (0, U0).
    """
    E_par, E_imp = (np.asarray(e, dtype=float) for e in energias)
    if len(E_par) < (n + 1) // 2 or len(E_imp) < n // 2:
        raise ValueError("No hay suficientes energías.")
    if np.any((E_par <= 0) | (E_par >= U0)) or np.any((E_imp <= 0) | (E_imp >= U0)):
        raise ValueError("Las energías deben cumplir 0 < E < U0.")

    x = np.linspace(-x_ext * a, x_ext * a, steps)
    dentro = np.abs(x) <= a
    psi = np.empty((n, steps))

    for j in range(n):
        i = j // 2
        E = E_par[i] if j % 2 == 0 else E_imp[i]
        k = np.sqrt(2 * M * E) / hbar
        q = np.sqrt(2 * M * (U0 - E)) / hbar
        A = np.sqrt(q / (a * q + 1))
        if j % 2 == 0:
            B = A * np.cos(k * a) * np.exp(q * a)
            psi[j] = np.where(dentro, A * np.cos(k * x), B * np.exp(-q * np.abs(x)))
        else:
            B = A * np.sin(k * a) * np.exp(q * a)
            psi[j] = np.where(dentro, A * np.sin(k * x),
                              np.sign(x) * B * np.exp(-q * np.abs(x)))

    return x, psi

def is_ortonormal(energias,M,hbar,U0,a,*,n=11):
    ''' 
    Function that shows a plot of S, the overlap matrix, that is used in quantum mechanics to measure 
    the orthogonality of the functions. The S_ij element is given by the internal product of psi_i and psi_j.

    Args:
        energias:   tuple that contains even and odd energies in arrays.
        M:          mass of the particle
        hbar:       value of Planck reduced constant
        U0:         value of the potential outside the well
        a:          width of the well
        n:          number of eigenstates to includei n matrix
    
    
    '''
    x, psi = eigenestados(energias,M,hbar,U0,a,n=n)
    S = np.trapezoid(psi[:, None, :] * psi[None, :, :], x, axis=-1)
    fig, ax = plt.subplots()
    im = ax.imshow(S, cmap="Greys", vmin=-1, vmax=1)
    fig.colorbar(im, ax=ax, label=r"$\langle\psi_i|\psi_j\rangle$")
    ax.set_xticks(range(len(S)), labels=range(1, len(S)+1))
    ax.set_yticks(range(len(S)), labels=range(1, len(S)+1))
    ax.set_xlabel("j")
    ax.set_ylabel("i")
    ax.set_title(f"Matriz de traslape $S$")
    plt.show()


# DATA:
hbar = 1
a = 1
M = 1/8 * np.pi**2
V0 = (hbar*np.pi)**2/(2*M*(2*a)**2)    
U0 = 110
mumax= np.sqrt(U0/V0)

# RESULTS:
energias = energias_permisibles(mumax)
plot_energies(energias)
plot_eigenstates(energias,M,hbar,U0,a)
is_ortonormal(energias,M,hbar,U0,a)
\end{lstlisting}

\section{Códigos de problema 1 Inciso 3(a)} \label{ap: sec1prob3a}

\begin{lstlisting}[language=Python, caption={Problema 1 Inciso 3(a).}]
from pathlib import Path
import numpy as np
from scipy.optimize import brentq
import matplotlib.pyplot as plt

# Compatibilidad: np.trapezoid (numpy >= 2.0) o np.trapz (versiones anteriores)
trapezoid = getattr(np, "trapezoid", None) or np.trapz

# Carpeta donde se guardan las figuras (miniproyecto1-cuantica/figuras)
FIG = Path(__file__).resolve().parent.parent / "figuras"
FIG.mkdir(exist_ok=True)

# =============================================================================
# DATA (mismos valores que el inciso 2)
# =============================================================================
hbar = 1
a = 1
M = 1/8 * np.pi**2
V0 = (hbar*np.pi)**2/(2*M*(2*a)**2)
U0 = 110
mumax = np.sqrt(U0/V0)


# =============================================================================
# Funciones del inciso 2
# =============================================================================
def energias_permisibles(mumax: float, *, eps=1e-12, a=1, hbar=1, M=1/8*np.pi**2):
    '''
    Calcula las energías permisibles del pozo finito simétrico resolviendo las
    ecuaciones (14) y (17). Los estados pares tienen mu en (2n, 2n+1) y los
    impares en (2n+1, 2n+2).

    Returns:
        Tuple(array, array) con las energías pares e impares respectivamente.
    '''
    f = lambda x: x*np.sin(np.pi*x/2) - np.sqrt(mumax**2 - x**2)*np.cos(np.pi*x/2)
    g = lambda x: np.sqrt(mumax**2 - x**2)*np.sin(np.pi*x/2) + x*np.cos(np.pi*x/2)
    raicespar, raicesimpar = [], []
    n = 0
    while 2*n < mumax:
        low, high = 2*n + eps, min(2*n + 1 - eps, mumax)
        if f(low)*f(high) < 0:
            raicespar.append(brentq(f, low, high))
        low += 1; high = min(high + 1, mumax)
        if low > mumax:
            break
        if g(low)*g(high) < 0:
            raicesimpar.append(brentq(g, low, high))
        n += 1
    E1 = (hbar*np.pi)**2/(2*M*(2*a)**2)
    return E1*np.array(raicespar)**2, E1*np.array(raicesimpar)**2


def eigenestados(energias, M, hbar, U0, a, *, n=8, x_ext=3, steps=6001):
    '''
    Evalúa los primeros n eigenestados normalizados (amplitudes analíticas)
    ordenados por energía creciente.

    Returns:
        x (steps,), psi (n, steps), E (n,)
    '''
    E_par, E_imp = energias
    x = np.linspace(-x_ext*a, x_ext*a, steps)
    dentro = np.abs(x) <= a
    psi = np.empty((n, steps))
    E = np.empty(n)
    for j in range(n):
        i = j // 2
        E[j] = E_par[i] if j % 2 == 0 else E_imp[i]
        k = np.sqrt(2*M*E[j])/hbar
        q = np.sqrt(2*M*(U0 - E[j]))/hbar
        A = np.sqrt(q/(a*q + 1))
        if j % 2 == 0:
            B = A*np.cos(k*a)*np.exp(q*a)
            psi[j] = np.where(dentro, A*np.cos(k*x), B*np.exp(-q*np.abs(x)))
        else:
            B = A*np.sin(k*a)*np.exp(q*a)
            psi[j] = np.where(dentro, A*np.sin(k*x), np.sign(x)*B*np.exp(-q*np.abs(x)))
    return x, psi, E


# =============================================================================
# Funciones del inciso 3
# =============================================================================
def evolucion(psi, cn, E, t, *, hbar=1):
    '''
    phi(x,t) = sum_n c_n phi_n(x) exp(-i E_n t / hbar)  (Ec. 2).

    Returns:
        Arreglo complejo de forma (len(t), len(x)).
    '''
    phi = np.zeros((len(t), psi.shape[1]), dtype=complex)
    for n in range(len(E)):
        if cn[n] != 0:
            phi += cn[n]*np.exp(-1j*E[n]*t/hbar)[:, None]*psi[n][None, :]
    return phi


def valor_esperado_x(x, phi):
    '''
    <x>(t) = int x |phi|^2 dx / int |phi|^2 dx
    '''
    rho = np.abs(phi)**2
    return trapezoid(x*rho, x, axis=1)/trapezoid(rho, x, axis=1)


def mapa_densidad(x, t, phi, titulo, archivo, *, a=1):
    '''
    Gráfica bidimensional de |phi(x,t)|^2.
    '''
    fig, ax = plt.subplots(figsize=(7, 4.5))
    im = ax.pcolormesh(x, t, np.abs(phi)**2, shading="auto", cmap="viridis")
    fig.colorbar(im, ax=ax, label=r"$|\phi(x,t)|^2$")
    ax.axvline(-a, color="red", lw=0.8, ls="--")
    ax.axvline(a, color="red", lw=0.8, ls="--")
    ax.set_xlabel(r"$x$")
    ax.set_ylabel(r"$t$")
    ax.set_title(titulo)
    fig.tight_layout()
    fig.savefig(archivo, dpi=200)
    plt.show()


def grafica_x(t, xm, titulo, archivo, *, a=1):
    '''
    Gráfica del valor esperado de la posición en función del tiempo.
    '''
    fig, ax = plt.subplots(figsize=(7, 3.5))
    ax.plot(t, xm, lw=1)
    ax.axhline(-a, color="red", lw=0.6, ls="--")
    ax.axhline(a, color="red", lw=0.6, ls="--")
    ax.set_xlabel(r"$t$")
    ax.set_ylabel(r"$\langle x\rangle$")
    ax.set_title(titulo)
    ax.grid(alpha=0.3)
    fig.tight_layout()
    fig.savefig(archivo, dpi=200)
    plt.show()


# =============================================================================
# RESULTS: inciso 3(a)
# =============================================================================
if __name__ == "__main__":
    energias = energias_permisibles(mumax)
    N = len(energias[0]) + len(energias[1])          # N = 11
    x, psi, E = eigenestados(energias, M, hbar, U0, a, n=N)

    # Coeficientes de la Ec. (4) y constante de normalización
    coef = np.array([1, 0.7, -1.2j, 2 - 1j])
    A = 1/np.sqrt(np.sum(np.abs(coef)**2))          # A = 1/sqrt(7.93)
    cn = np.zeros(N, dtype=complex)
    cn[:4] = A*coef

    t = np.linspace(0, 10, 1201)
    phi = evolucion(psi, cn, E, t)

    print(f"A = {A:.4f}")
    print(f"Norma en t=0: {trapezoid(np.abs(phi[0])**2, x):.10f}")

    mapa_densidad(x, t, phi, r"Inciso 3(a): $|\phi(x,t)|^2$", FIG / "p1_3a_mapa.png")
    grafica_x(t, valor_esperado_x(x, phi), r"Inciso 3(a): $\langle x\rangle(t)$", FIG / "p1_3a_x.png")
\end{lstlisting}

\section{Códigos de problema 1 Inciso 3(b)} \label{ap: sec1prob3b}

\begin{lstlisting}[language=Python, caption={Problema 1 Inciso 3(b).}]
from pathlib import Path
import numpy as np
from scipy.optimize import brentq
import matplotlib.pyplot as plt
from matplotlib.animation import FuncAnimation, FFMpegWriter

trapezoid = getattr(np, "trapezoid", None) or np.trapz

FIG = Path(__file__).resolve().parent.parent / "figuras"
FIG.mkdir(exist_ok=True)
VID = Path(__file__).resolve().parent.parent / "videos"
VID.mkdir(exist_ok=True)

try:
    import imageio_ffmpeg
    plt.rcParams["animation.ffmpeg_path"] = imageio_ffmpeg.get_ffmpeg_exe()
except ImportError:
    pass

hbar = 1
a = 1
M = 1/8 * np.pi**2
V0 = (hbar*np.pi)**2/(2*M*(2*a)**2)
U0 = 110
mumax = np.sqrt(U0/V0)


def energias_permisibles(mumax: float, *, eps=1e-12, a=1, hbar=1, M=1/8*np.pi**2):
    '''
    Calcula las energías permisibles del pozo finito simétrico resolviendo las
    ecuaciones (14) y (17).

    Returns:
        Tuple(array, array) con las energías pares e impares respectivamente.
    '''
    f = lambda x: x*np.sin(np.pi*x/2) - np.sqrt(mumax**2 - x**2)*np.cos(np.pi*x/2)
    g = lambda x: np.sqrt(mumax**2 - x**2)*np.sin(np.pi*x/2) + x*np.cos(np.pi*x/2)
    raicespar, raicesimpar = [], []
    n = 0
    while 2*n < mumax:
        low, high = 2*n + eps, min(2*n + 1 - eps, mumax)
        if f(low)*f(high) < 0:
            raicespar.append(brentq(f, low, high))
        low += 1; high = min(high + 1, mumax)
        if low > mumax:
            break
        if g(low)*g(high) < 0:
            raicesimpar.append(brentq(g, low, high))
        n += 1
    E1 = (hbar*np.pi)**2/(2*M*(2*a)**2)
    return E1*np.array(raicespar)**2, E1*np.array(raicesimpar)**2


def eigenestados(energias, M, hbar, U0, a, *, n=8, x_ext=3, steps=6001):
    '''
    Evalúa los primeros n eigenestados normalizados ordenados por energía.

    Returns:
        x (steps,), psi (n, steps), E (n,)
    '''
    E_par, E_imp = energias
    x = np.linspace(-x_ext*a, x_ext*a, steps)
    dentro = np.abs(x) <= a
    psi = np.empty((n, steps))
    E = np.empty(n)
    for j in range(n):
        i = j // 2
        E[j] = E_par[i] if j % 2 == 0 else E_imp[i]
        k = np.sqrt(2*M*E[j])/hbar
        q = np.sqrt(2*M*(U0 - E[j]))/hbar
        A = np.sqrt(q/(a*q + 1))
        if j % 2 == 0:
            B = A*np.cos(k*a)*np.exp(q*a)
            psi[j] = np.where(dentro, A*np.cos(k*x), B*np.exp(-q*np.abs(x)))
        else:
            B = A*np.sin(k*a)*np.exp(q*a)
            psi[j] = np.where(dentro, A*np.sin(k*x), np.sign(x)*B*np.exp(-q*np.abs(x)))
    return x, psi, E


def evolucion(psi, cn, E, t, *, hbar=1):
    '''
    phi(x,t) = sum_n c_n phi_n(x) exp(-i E_n t / hbar)  (Ec. 2).
    '''
    phi = np.zeros((len(t), psi.shape[1]), dtype=complex)
    for n in range(len(E)):
        if cn[n] != 0:
            phi += cn[n]*np.exp(-1j*E[n]*t/hbar)[:, None]*psi[n][None, :]
    return phi


def valor_esperado_x(x, phi):
    '''
    <x>(t) = int x |phi|^2 dx / int |phi|^2 dx
    '''
    rho = np.abs(phi)**2
    return trapezoid(x*rho, x, axis=1)/trapezoid(rho, x, axis=1)


def coeficientes(x, psi, f):
    '''
    c_n = int phi_n^*(x) f(x) dx  (Ec. 3) con la regla del trapecio.
    '''
    return np.array([trapezoid(np.conj(p)*f, x) for p in psi])


def gaussiano(x, b, c, v0, *, M=1/8*np.pi**2, hbar=1):
    '''
    Pulso Gaussiano de la Ec. (5) con A = (2/(pi b^2))^(1/4).
    '''
    A = (2/(np.pi*b**2))**0.25
    return A*np.exp(-(x - c)**2/b**2)*np.exp(1j*M*v0*x/hbar)


def mapa_densidad(x, t, phi, titulo, archivo, *, a=1):
    '''
    Gráfica bidimensional de |phi(x,t)|^2.
    '''
    fig, ax = plt.subplots(figsize=(7, 4.5))
    im = ax.pcolormesh(x, t, np.abs(phi)**2, shading="auto", cmap="viridis")
    fig.colorbar(im, ax=ax, label=r"$|\phi(x,t)|^2$")
    ax.axvline(-a, color="red", lw=0.8, ls="--")
    ax.axvline(a, color="red", lw=0.8, ls="--")
    ax.set_xlabel(r"$x$")
    ax.set_ylabel(r"$t$")
    ax.set_title(titulo)
    fig.tight_layout()
    fig.savefig(archivo, dpi=200)
    plt.show()


def grafica_x(t, xm, titulo, archivo, *, a=1):
    '''
    Gráfica del valor esperado de la posición en función del tiempo.
    '''
    fig, ax = plt.subplots(figsize=(7, 3.5))
    ax.plot(t, xm, lw=1)
    ax.axhline(-a, color="red", lw=0.6, ls="--")
    ax.axhline(a, color="red", lw=0.6, ls="--")
    ax.set_xlabel(r"$t$")
    ax.set_ylabel(r"$\langle x\rangle$")
    ax.set_title(titulo)
    ax.grid(alpha=0.3)
    fig.tight_layout()
    fig.savefig(archivo, dpi=200)
    plt.show()


def grafica_fg(x, f, g, archivo, *, a=1):
    '''
    Compara el pulso inicial f(x) con su reconstrucción g(x) usando los N estados.
    '''
    fig, ax = plt.subplots(figsize=(7, 3.5))
    ax.plot(x, np.abs(f)**2, lw=1, label=r"$|f(x)|^2$")
    ax.plot(x, np.abs(g)**2, lw=1, ls="--", label=r"$|g(x)|^2$")
    ax.axvline(-a, color="red", lw=0.6, ls="--")
    ax.axvline(a, color="red", lw=0.6, ls="--")
    ax.set_xlim(-2*a, 2*a)
    ax.set_xlabel(r"$x$")
    ax.set_title(r"Pulso inicial $f(x)$ y su reconstrucción $g(x)$")
    ax.legend()
    ax.grid(alpha=0.3)
    fig.tight_layout()
    fig.savefig(archivo, dpi=200)
    plt.show()


def video(x, t, phi, archivo, *, a=1, U0=110, fps=30, titulo=""):
    '''
    Video mp4 de la evolución de |phi(x,t)|^2 junto con el potencial.
    '''
    rho = np.abs(phi)**2
    ymax = 1.1*rho.max()
    fig, ax = plt.subplots(figsize=(8, 4.5))
    pot = np.where(np.abs(x) <= a, 0, U0)/U0*ymax*0.9
    ax.fill_between(x, 0, pot, color="red", alpha=0.12, label=r"$U(x)$ (escalado)")
    linea, = ax.plot(x, rho[0], lw=1.2, label=r"$|\phi(x,t)|^2$")
    re, = ax.plot(x, np.real(phi[0]), lw=0.6, alpha=0.6, label=r"Re$\,\phi$")
    ax.set_xlim(x[0], x[-1])
    ax.set_ylim(-np.sqrt(ymax), ymax)
    ax.set_xlabel(r"$x$")
    ax.legend(loc="upper right")
    ax.grid(alpha=0.3)
    texto = ax.set_title("")

    def update(i):
        linea.set_ydata(rho[i])
        re.set_ydata(np.real(phi[i]))
        texto.set_text(f"{titulo}   t = {t[i]:.3f}")
        return linea, re, texto

    anim = FuncAnimation(fig, update, frames=len(t), blit=False)
    anim.save(archivo, writer=FFMpegWriter(fps=fps, bitrate=2400))
    plt.close(fig)


if __name__ == "__main__":
    energias = energias_permisibles(mumax)
    N = len(energias[0]) + len(energias[1])
    x, psi, E = eigenestados(energias, M, hbar, U0, a, n=N)

    b, c, v0 = 0.2, 0.4, 8
    f = gaussiano(x, b, c, v0)
    cn = coeficientes(x, psi, f)
    g = sum(cn[n]*psi[n] for n in range(N))
    P = np.sum(np.abs(cn)**2)

    print(f"A = {(2/(np.pi*b**2))**0.25:.4f}")
    print(f"Norma de f: {trapezoid(np.abs(f)**2, x):.6f}")
    print("|c_n|^2 =", np.round(np.abs(cn)**2, 4))
    print(f"Suma |c_n|^2 = {P:.4f}")
    print(f"<E> de g = {np.sum(np.abs(cn)**2*E)/P:.2f}")

    grafica_fg(x, f, g, FIG / "p1_3b_fg.png")

    t = np.linspace(0, 3, 1201)
    phi = evolucion(psi, cn, E, t)
    mapa_densidad(x, t, phi, r"Inciso 3(b): $|\phi(x,t)|^2$ del pulso Gaussiano", FIG / "p1_3b_mapa.png")
    grafica_x(t, valor_esperado_x(x, phi), r"Inciso 3(b): $\langle x\rangle(t)$ del pulso Gaussiano", FIG / "p1_3b_x.png")

    print("Generando video (tarda alrededor de un minuto)...")
    t_v = np.linspace(0, 3, 451)
    sel = np.abs(x) <= 2*a
    video(x[sel], t_v, evolucion(psi[:, sel], cn, E, t_v), VID / "pulso_pozo_simetrico.mp4",
          titulo="Pozo simétrico")
    print("Video guardado en videos/pulso_pozo_simetrico.mp4")
\end{lstlisting}

\section{Códigos de problema 1 Inciso 4(b)} \label{ap: sec1prob4b}

\begin{lstlisting}[language=Python, caption={Problema 1 Inciso 4(b).}]
import numpy as np
from scipy.optimize import brentq

hbar = 1
a = 1
M = 1/8 * np.pi**2
V0 = (hbar*np.pi)**2/(2*M*(2*a)**2)
U0 = 110
mumax = np.sqrt(U0/V0)


def energias_permisibles(mumax: float, *, eps=1e-12, a=1, hbar=1, M=1/8*np.pi**2):
    '''
    Calcula las energías permisibles del pozo finito simétrico resolviendo las
    ecuaciones (14) y (17).

    Returns:
        Tuple(array, array) con las energías pares e impares respectivamente.
    '''
    f = lambda x: x*np.sin(np.pi*x/2) - np.sqrt(mumax**2 - x**2)*np.cos(np.pi*x/2)
    g = lambda x: np.sqrt(mumax**2 - x**2)*np.sin(np.pi*x/2) + x*np.cos(np.pi*x/2)
    raicespar, raicesimpar = [], []
    n = 0
    while 2*n < mumax:
        low, high = 2*n + eps, min(2*n + 1 - eps, mumax)
        if f(low)*f(high) < 0:
            raicespar.append(brentq(f, low, high))
        low += 1; high = min(high + 1, mumax)
        if low > mumax:
            break
        if g(low)*g(high) < 0:
            raicesimpar.append(brentq(g, low, high))
        n += 1
    E1 = (hbar*np.pi)**2/(2*M*(2*a)**2)
    return E1*np.array(raicespar)**2, E1*np.array(raicesimpar)**2


def energias_unificadas(mumax: float, *, V0=1):
    '''
    Resuelve la ecuación característica unificada
        mu = n - (2/pi) arcsin(mu/mumax),   n = 1, 2, ...
    Hay una raíz por cada n mientras F(mumax) = mumax + 1 > n.

    Returns:
        Array con las energías E_n = V0 mu_n^2 en orden creciente.
    '''
    F = lambda mu, n: mu + (2/np.pi)*np.arcsin(mu/mumax) - n
    mus = []
    n = 1
    while F(mumax, n) > 0:
        mus.append(brentq(F, 0, mumax, args=(n,), xtol=1e-14))
        n += 1
    return V0*np.array(mus)**2


if __name__ == "__main__":
    E_par, E_imp = energias_permisibles(mumax)
    E_paridad = np.sort(np.concatenate((E_par, E_imp)))
    E_unif = energias_unificadas(mumax, V0=V0)

    print(f"N (paridad) = {len(E_paridad)},  N (unificada) = {len(E_unif)},  ceil(mumax) = {int(np.ceil(mumax))}")
    print(f"{'n':>3} {'Ecs. 14 y 17':>14} {'Unificada':>12} {'Diferencia':>12}")
    for n, (e1, e2) in enumerate(zip(E_paridad, E_unif), start=1):
        print(f"{n:>3} {e1:>14.4f} {e2:>12.4f} {abs(e1 - e2):>12.1e}")
    print(f"Diferencia máxima: {np.max(np.abs(E_paridad - E_unif)):.1e}")
\end{lstlisting}


\section{Códigos de problema 2 Inciso 6} \label{ap: sec2prob6}
Se utilizó Julia y se usó la IA mayormente para resolver dudas y graficar de mejor manera. Todo el código y lógica fue hecho por los autores y estos se hacen responsables de su veracidad.
\begin{lstlisting}[language=Julia, caption={Problema 2 Inciso 6}]
h = 1
L = 2
M = (pi^2)/8
U1 = 85
c = 2
U2 = c*U1
E = range(0.001, 84.999, length=3000)

Tangente = tan.((L/h)*sqrt.(2*M*E))
F2 = (sqrt.(E) .* (sqrt.(U1 .- E) .+ sqrt.(U2 .- E))) ./ (E .- sqrt.((U1 .- E) .* (U2 .- E)))

for i in 1:(length(Tangente)-1) #quiero quitar los valores que se rompen en la tangente
    if abs(Tangente[i+1] - Tangente[i]) > 160
        Tangente[i] = NaN
    end
end
for i in 1:(length(F2)-1) #quitamos lineas horizontales, ya que no representan algo
    if abs(F2[i+1] - F2[i]) > 20
        F2[i] = NaN
    end
end
g(e) = (L/h)*sqrt(2*M*e) - atan(sqrt(U1-e)/sqrt(e)) - atan(sqrt(U2-e)/sqrt(e)) #tomamos la fecuacion de 2.8
E_puntos = Float64[]
for n in 1:9 #encontramos que puntos corresponden
    f(e) = g(e) - (n - 1) * pi
    raiz = find_zero(f, (0.001, 84.999))
    push!(E_puntos, raiz)
end
Y_puntos = tan.((L/h) .* sqrt.(2*M .* E_puntos))# Evaluamos
grafica = plot(E, Tangente, ylims=(-160,160), label="tan(ωL)", xlabel="Energía E", ylabel="funciones(E)", title="Grafico de eigenenergias")
plot!(grafica,E,F2)
scatter!(E_puntos, Y_puntos, color=:red, markersize=5, label="Eigenenergías corte de funciones")
offsets_y = [25.0, -25.0, 70.0, -65.0, 25.0, -25.0, 25.0, 0.0, 25.0]

for i in 1:9
    e_val = round(E_puntos[i], digits=1)
    
    if Y_puntos[i] < -100
        annotate!(grafica, E_puntos[i] + 2.5, Y_puntos[i], text("E$i = $e_val", 7, :left, :bold))
    else
        annotate!(grafica, E_puntos[i], Y_puntos[i] + offsets_y[i], text("E$i = $e_val", 7, :center, :bold))
    end
end

display(grafica)
x = range(-1.5, L + 1.5, length=1000) #hagamos nuestros valores de x a graficar
region1 = x .< 0.0
region2 = (x .>= 0.0) .& (x .<= L)
region3 = x .> L
function phi_n(x, E_n, region1, region2, region3; h=1.0, L=2.0, M=(pi^2)/8.0, U1=85.0, U2=170.0)
    wn = sqrt(2 * M * E_n) / h
    k  = sqrt(2 * M * (U1 - E_n)) / h
    y  = sqrt(2 * M * (U2 - E_n)) / h
    beta = (sin(wn * L) * (k^2 + wn^2)) / (wn * (y + k))
    Calrevez = (1 / (2 * k)) + (L / 2) * (1 + (k^2 / wn^2)) + (sin(2 * wn * L) / (4 * wn)) * (1 - (k^2 / wn^2)) + (k / wn^2) * (sin(wn * L)^2) + (beta^2 / (2 * y))     
    C = 1 / sqrt(Calrevez)
    phi = zeros(length(x))
    phi[region1] .= C .* exp.(k .* x[region1])
    phi[region2] .= C .* (cos.(wn .* x[region2]) .+ (k / wn) .* sin.(wn .* x[region2]))
    phi[region3] .= C .* beta .* exp.(-y .* (x[region3] .- L))
    return phi
end
U = zeros(length(x))
U[region1] .= U1
U[region2] .= 0
U[region3] .= U2
graficas = plot(x,U,color=:black,label="U(x)",xlabel = "x",ylabel = "Energia",title = "Grafica de eigenestados",size(850,600))
for (n, E_n) in enumerate(E_puntos)
    phi_vec =phi_n(x, E_n, region1, region2, region3)
    plot!(graficas, x, E_n .+ 3* phi_vec, label="n = $n")
end
display(graficas)
b = 0.2
cancho = 0.4
v0 = 0.0
A = (2.0 / (pi * b^2))^(0.25)

phi_0 = A .* exp.(-((x .- cancho).^2) ./ (b^2)) .* exp.(1im .* (M * v0 / h) .* x)
dx = x[2] - x[1]
N_estados = length(E_puntos)
c_n = zeros(ComplexF64,N_estados)
for n in 1:length(E_puntos)
    phi_vec = phi_n(x, E_puntos[n], region1, region2, region3)
    c_n[n] = sum(phi_vec.*phi_0)*dx
end
for n in 1:length(c_n)
    println("c_$n = ", real.(c_n[n]))
end
suma = 0 
for or in 1:length(c_n)
    suma += (c_n[or]).^2
end
real(suma)
t_final = 5.0
t = range(0.0, t_final, length=200)
Nx = length(x)
Nt = length(t)
P_xt = zeros(Float64, Nx, Nt)
phi_m = zeros(Float64, Nx, N_estados)
for n in 1:N_estados
    phi_m[:, n] = phi_n(x, E_puntos[n], region1, region2, region3)
end
for j in 1:Nt
    tj = t[j]
    for k in 1:Nx
        xk = x[k]
        sumas = 0.0 + 0.0im
        
        for n in 1:N_estados
            phi_n_x = phi_m[k, n]
            cn = c_n[n]
            En = E_puntos[n]
            
            for m in 1:N_estados
                phi_m_x = phi_m[k, m]
                cm_conj = conj(c_n[m])
                Em = E_puntos[m]
                prob = cn * cm_conj * phi_n_x * phi_m_x * exp(-1im * (En - Em) * tj / h)
                sumas += prob
            end
        end
        P_xt[k, j] = real(sumas)
    end
end
grafica_2D = heatmap(x, t, P_xt',xlabel = "Posición x", ylabel = "Tiempo t",title = "Evolución de g(x,t)^2",color = :viridis,size = (800, 500))

display(grafica_2D)
valor_esperado = zeros(Float64, Nt)
for j in 1:Nt
    P_columna = P_xt[:, j]
    valor_esperado[j] = sum(x .* P_columna) * dx
end
valoresperadografica = plot(t, valor_esperado,linewidth = 2,color = :blue,xlabel = "Tiempo t",ylabel = "<x>(t)",title = "Valor esperado de x en el tiempo",size = (800, 400))
display(valoresperadografica)

Tope = maximum(P_xt)
animacion = @animate for j in 1:Nt
    tj = t[j]
    P_actual = P_xt[:, j]
    plot(
        x, P_actual,
        linewidth = 2,
        color = :blue,  label = "P(x,t) = |g(x,t)|²",xlabel = "x",ylabel = "P(x,t)",title = "Evolución del Pulso (t = $(round(tj, digits=2)))",xlims = (-0.3, L + 0.3),ylims = (0.0, Tope * 1.05),legend = :topright,size = (800, 450))
    vline!([0.0], color = :black, linestyle = :solid, linewidth = 1.5, label = "Paredes")
    vline!([L], color = :black, linestyle = :solid, linewidth = 1.5, label = "")
end
mp4(animacion, "pulso.mp4", fps = 15)
\end{lstlisting}